\chapter{Constructing the flight model step by step}
\section{Start from a free-body diagram, not a guessed transfer function}
Our first modelling choice is a system boundary: the complete flying aircraft,
including battery, frame, motors, electronics and propellers. Its total mass is
$m=\sum_i m_i$, and its centre of mass is
$r_C=\sum_i m_i r_i/m$. Choosing the body origin at $r_C$ removes the extra
moment caused by expressing angular momentum about an accelerating offset point.
Neither a battery mass nor a frame mass is the complete $m$ in Newton's law.

Choose a local world frame with upward unit vector $e_3$ and a forward-left-up
body frame. Let the four thrust forces be $T_i e_3$ in body coordinates.
There are three categories of external force: rotated rotor thrust, weight,
and aerodynamic/disturbance force. Internal frame forces cancel when summing
over the system boundary. Thus the force balance is constructed as
\begin{align}
\sum f_W&=\underbrace{R e_3(T_1+T_2+T_3+T_4)}_{\text{rotors}}
 -\underbrace{mg e_3}_{\text{weight}}
 -\underbrace{d_vv_W}_{\text{linear drag}}+f_w,\\
\dot p_W&=v_W,\qquad m\dot v_W=\sum f_W.
\end{align}
Dividing by $m$ produces acceleration, not velocity. Integrating acceleration
once gives velocity and a second time gives position. This is why a translational
position transfer function has two low-frequency integrations before damping,
control or other dynamics are introduced.

For each motor at $r_i=(x_i,y_i,0)$, the torque about the centre of mass is
$(y_iT_i,-x_iT_i,0)$. Add rotor drag-reaction torques with alternating signs.
The sum is the allocation equation derived in the mechanics chapter. Newton's
angular counterpart is $J\dot\omega+\omega\times J\omega=\tau$. We keep the
cross term in the nonlinear plant and explicitly neglect its second-order
perturbation when constructing the hover linear model.

\section{Write the nonlinear differential equations before linearizing}
Write $T=\sum T_i$ and use the third column of $R$. The translational equations
are, component by component,
\begin{align}
m\ddot x&=T(c_\psi s_\theta c_\phi+s_\psi s_\phi)-d_v\dot x+f_{wx},\\
m\ddot y&=T(s_\psi s_\theta c_\phi-c_\psi s_\phi)-d_v\dot y+f_{wy},\\
m\ddot h&=Tc_\theta c_\phi-mg-d_v\dot h+f_{wz}.
\end{align}
For diagonal inertia,
\begin{align}
\dot p&=[\tau_x+(J_y-J_z)qr]/J_x,\\
\dot q&=[\tau_y+(J_z-J_x)rp]/J_y,\\
\dot r&=[\tau_z+(J_x-J_y)pq]/J_z.
\end{align}
These six acceleration equations need kinematics as well: $\dot p_W=v_W$
and quaternion attitude propagation. The state is position, velocity, attitude
and angular velocity; four additional motor-lag states make the actuator
dynamics explicit. Quaternion norm is constrained to one, so four quaternion
coordinates represent three independent orientation degrees of freedom.

\section{Find the hover equilibrium by substitution}
An equilibrium has zero time derivatives. Choose $v_0=0$, $\omega_0=0$,
$\phi_0=\theta_0=0$ and any fixed heading $\psi_0$. Substitution in the
vertical equation gives $0=T_0-mg$, so $T_0=mg$. Zero torque with a symmetric
allocation gives $T_{i0}=mg/4$. If normalized desired thrust satisfies
$T_i=T_{\max}u_i$, then
\begin{equation}
u_{i0}=c_0=\frac{mg}{4T_{\max}}.
\end{equation}
For $m=1.2$ kg and $T_{\max}=6$ N, each motor must supply $2.941995$ N and
$c_0=0.4903325$. This is a model calculation, not an identified ESC duty cycle.
The thrust-to-weight ratio of this assumed model is $24/11.76798\simeq2.0394$.
It is not evidence that the real 1400KV/9047 combination has that reserve.

\section{Linearize by explicitly retaining first-order terms}
Define perturbations $\delta T=T-T_0$, $\delta\phi=\phi-\phi_0$ and so on.
The first-order Taylor expansion of $\sin\epsilon$ is $\epsilon$; that of
$\cos\epsilon$ is $1$. Products such as $\delta T\delta\theta$ contain two
small perturbations and are second order. At heading zero,
\begin{align}
(T_0+\delta T)\sin\delta\theta\cos\delta\phi
 &\simeq T_0\delta\theta,\\
-(T_0+\delta T)\sin\delta\phi&\simeq-T_0\delta\phi,\\
(T_0+\delta T)\cos\delta\theta\cos\delta\phi-mg&\simeq\delta T.
\end{align}
Using $T_0/m=g$, define $\gamma=d_v/m=0.1$ s$^{-1}$ and obtain
\begin{align}
\delta\ddot x+\gamma\,\delta\dot x&=g\,\delta\theta+f_{wx}/m,\\
\delta\ddot y+\gamma\,\delta\dot y&=-g\,\delta\phi+f_{wy}/m,\\
\delta\ddot h+\gamma\,\delta\dot h&=\delta T/m+f_{wz}/m.
\end{align}
Here $\gamma$ is the damping coefficient and $\delta x$ denotes a perturbation.
Below we omit the perturbation prefix to keep the equations legible. At a nonzero heading, rotate the two
horizontal acceleration equations by the heading angle. This is a local
linearization, not a claim that Euler coordinates work globally.

The body-rate products $qr,rp,pq$ vanish to first order about zero rates.
Near level attitude $\dot\phi\simeq p$, $\dot\theta\simeq q$, $\dot\psi\simeq r$.
Thus $J_x\ddot\phi=\tau_x$, $J_y\ddot\theta=\tau_y$ and
$J_z\dot r=\tau_z$. There is intentionally no yaw-angle feedback requirement.

\section{Build each open-loop transfer function from the differential equation}
A transfer function describes a linear input/output relation with zero initial
perturbations. For $\dot y=f$, the unilateral Laplace transform is
$sY(s)-y(0)=F(s)$. For $\ddot y$, it is
$s^2Y(s)-sy(0)-\dot y(0)$. Setting the initial perturbations to zero removes
those initial-condition terms; it does not mean the actual aircraft has no
initial state. Initial-condition responses must be added separately.

From $J_x\ddot\phi=\tau_x$, first transform to
$J_xs^2\Phi=\mathcal{T}_x$, then divide both sides by $J_xs^2\mathcal{T}_x$:
\begin{equation}
\frac{\Phi(s)}{\mathcal{T}_x(s)}=\frac{1}{J_xs^2},\quad
\frac{P(s)}{\mathcal{T}_x(s)}=\frac{1}{J_xs}.
\end{equation}
Use the allocation $\tau_x=4aT_{\max}u_\phi$ and the motor lag
$\tau_m\dot u_{\rm act}+u_{\rm act}=u_{\rm cmd}$.
Transforming the latter gives
$(\tau_ms+1)U_{\rm act}=U_{\rm cmd}$. Therefore the cascade is
\begin{equation}
G_{p}(s)=\frac{P_d(s)}{U_\phi(s)}=
\frac{b_x}{s(\tau_ms+1)},\qquad
G_{\phi}(s)=\frac{\Phi_d(s)}{U_\phi(s)}=
\frac{b_x}{s^2(\tau_ms+1)},
\end{equation}
where $b_x=(4aT_{\max}/J_x)(180/\pi)$ because the controller uses degrees.
Pitch has $b_y$ instead; yaw rate has $b_\psi=4\kappa T_{\max}(180/\pi)/J_z$.
The rate plant is second order with lag and the angle plant third order.
Without actuator lag these reduce to first and second order respectively.

For altitude, $\ddot h+\gamma\dot h=b_hu_{\rm act}$ with
$b_h=4T_{\max}/m=20$ m/s$^2$/command. Transforming each derivative gives
\begin{equation}
G_h(s)=\frac{H(s)}{U_h(s)}=
\frac{b_h}{s(s+\gamma)(\tau_ms+1)}.
\end{equation}
This retains drag and actuator lag. The often-used double integrator $b_h/s^2$
is obtained only after explicitly neglecting both. Horizontal position from
small angle is $X/\Theta=g/[s(s+\gamma)]$ and
$Y/\Phi=-g/[s(s+\gamma)]$ in radians. Degree-valued angle inputs require an
additional factor $\pi/180$.

\section{Relate normalized thrust to a possible rotor-speed model}
A better identified actuator could use $T=k_T\Omega^2$ and
$Q=k_Q\Omega^2$. Around nonzero hover speed $\Omega_0$, expand
\begin{equation}
k_T(\Omega_0+\Delta\Omega)^2
=k_T\Omega_0^2+2k_T\Omega_0\Delta\Omega+k_T(\Delta\Omega)^2.
\end{equation}
Discarding the quadratic perturbation gives
$\Delta T=2k_T\Omega_0\Delta\Omega$. A first-order identified command-to-speed
gain $k_\Omega/(\tau_\Omega s+1)$ then yields
$\Delta T/\Delta u=2k_T\Omega_0k_\Omega/(\tau_\Omega s+1)$.
The current normalized-thrust model replaces that unknown product with
$T_{\max}$. KV alone supplies none of $k_T,k_Q,k_\Omega,\tau_\Omega$ under
load. A higher-order model is useful only when its extra parameters are measured.

\chapter{Closed-loop transfer functions and stability calculations}
\section{Rate PI: derive the denominator without skipping the feedback algebra}
Let $G=b/[s(\tau_ms+1)]$, $C=K_p+K_i/s$ and $e=r_c-r$.
Then $r=GCe=GC(r_c-r)$. Move the feedback term to the left:
$(1+GC)r=GC r_c$. Divide by $1+GC$ and clear denominators:
\begin{equation}
T_r(s)=\frac{r}{r_c}=
\frac{b(K_ps+K_i)}{\tau_ms^3+s^2+bK_ps+bK_i}.
\end{equation}
The numerator zero at $s=-K_i/K_p$ changes the step response; denominator
poles alone do not completely describe overshoot. With $\tau_m=0$, matching
$s^2+2\zeta\omega_ns+\omega_n^2$ gives $K_p=2\zeta\omega_n/b$ and
$K_i=\omega_n^2/b$. This is an initial design followed by checking the cubic,
not an exact second-order description of the complete plant.

\section{Roll/pitch angle cascade: retain the inner loop dynamics}
Let $k_a=3$ s$^{-1}$ and $p_c=k_a(\phi_c-\phi)$. Since $p=s\phi$,
\begin{align}
\phi&=\frac{T_r}{s}k_a(\phi_c-\phi),\\
\left(1+\frac{k_aT_r}{s}\right)\phi&=\frac{k_aT_r}{s}\phi_c,\\
T_\phi(s)&=\frac{k_aT_r}{s+k_aT_r}.
\end{align}
Substitute the rate-loop numerator and denominator and collect powers:
\begin{equation}
\boxed{T_\phi(s)=\frac{bk_a(K_ps+K_i)}{
\tau_ms^4+s^3+bK_ps^2+b(K_i+k_aK_p)s+bk_aK_i}.}
\end{equation}
Pitch is identical only under the assumed $J_x=J_y$. This fourth-order
expression is more complete than replacing the inner rate loop by unity.
The ideal outer approximation $k_a/(s+k_a)$ is justified only well below the
inner closed-loop bandwidth, away from saturation and estimator error.

\section{Altitude PID and the distinction between height and velocity requests}
At fixed equilibrium, $e_h=h_c-h$, $\dot I_h=e_h$ and
$u=K_pe_h+K_iI_h+K_d(v_c-v)$, where the internally integrated target satisfies
$\dot h_c=v_c$. For differentiable consistent references,
$V_c=sH_c$, $V=sH$ and $I_h=(H_c-H)/s$, so
\begin{equation}
U=(K_p+K_i/s+K_ds)(H_c-H).
\end{equation}
Combine this with $G_h$ above and multiply the characteristic equation by
$s^2(s+\gamma)(\tau_ms+1)$:
\begin{align}
D_h(s)&=s^2(s+\gamma)(\tau_ms+1)
+b_h(K_ds^2+K_ps+K_i)\\
&=\tau_ms^4+(1+\gamma\tau_m)s^3
+(\gamma+b_hK_d)s^2+b_hK_ps+b_hK_i,\\
\boxed{T_h(s)}&=\boxed{\frac{b_h(K_ds^2+K_ps+K_i)}{D_h(s)}}.
\end{align}
The transfer from commanded vertical velocity to height is $T_h(s)/s$;
the transfer to vertical velocity is $T_h(s)$, with zero initial perturbations.
A discontinuous height step accompanied by an impulsive derivative reference
is not the actual stick interface. The mission generates bounded $v_c$; the
height target integrates it. If one instead steps $h_c$ while forcing $v_c=0$,
the numerator is $b_h(K_ps+K_i)$: the denominator is the same, but the reference
response differs. This prevents applying a textbook PID step response to the
wrong command interface.

When lag is neglected but drag retained, matching
$(s+p_h)(s^2+2\zeta_h\omega_hs+\omega_h^2)$ gives
\begin{align}
K_d&=(2\zeta_h\omega_h+p_h-\gamma)/b_h,\\
K_p&=(\omega_h^2+2p_h\zeta_h\omega_h)/b_h,\\
K_i&=p_h\omega_h^2/b_h.
\end{align}
The existing baseline was designed with $\gamma=0$, giving $K_d=0.175$.
Including the assumed drag gives $K_d=0.170$. The selected $\lambda=1.3$
candidate similarly gives $0.2275$ versus $0.2225$. These differences are shown
as model alternatives; gains are not silently changed without a new route test.

\section{What second-order parameters actually mean}
For a complex pair $s_{1,2}=-\sigma\pm j\omega_d$ with $\sigma>0$,
$\omega_n=\sqrt{\sigma^2+\omega_d^2}$ and $\zeta=\sigma/\omega_n$.
For the ideal normalized second-order low-pass without zeros,
\begin{equation}
M_p=e^{-\pi\zeta/\sqrt{1-\zeta^2}},\quad
t_p=\pi/\omega_d,\quad t_{s,2\%}\simeq4/(\zeta\omega_n).
\end{equation}
These follow from the decaying sinusoidal homogeneous response and its extrema;
the settling approximation bounds its exponential envelope approximately.
For example the ideal rate design $\zeta=0.9$, $\omega_n=12$ has
$\omega_d=5.2307$ rad/s and envelope settling about $0.3704$ s. Its ideal
no-zero overshoot is about $0.152\%$. The actual PI closed loop has a zero,
and actuator lag adds another pole. Therefore these figures describe the
chosen target polynomial, not a guaranteed measured aircraft step response.
The companion numerical analysis computes actual step responses for the
complete local transfer functions, including their zeros.

\section{Routh--Hurwitz: construct and evaluate the stability conditions}
For $a_3s^3+a_2s^2+a_1s+a_0$ with $a_3>0$, the Routh first column is
\begin{equation}
a_3,\quad a_2,\quad\frac{a_2a_1-a_3a_0}{a_2},\quad a_0.
\end{equation}
All entries must be strictly positive for strict left-half-plane roots, assuming
no degenerate zero rows. Applying this to the rate loop gives positive gains
and $K_p>\tau_mK_i$. Substituting the pole-matched PI gives
$\omega_n<2\zeta/\tau_m$. For $\tau_m=0.04$ s the bound is 45 rad/s;
the selected rate frequency 15.6 rad/s meets this local cubic condition.

For $a_4s^4+a_3s^3+a_2s^2+a_1s+a_0$, construct
\begin{equation}
\begin{array}{c|ccc}
s^4&a_4&a_2&a_0\\s^3&a_3&a_1&0\\
s^2&\Delta_2/a_3&a_0&0\\
s^1&\Delta_3/\Delta_2&0&0\\s^0&a_0&0&0
\end{array},
\end{equation}
where $\Delta_2=a_3a_2-a_4a_1$ and
$\Delta_3=a_1\Delta_2-a_3^2a_0$. With $a_4,a_3,a_0>0$, strict stability
requires $\Delta_2>0$ and $\Delta_3>0$. Positive polynomial coefficients alone
are insufficient. Substitution for the full angle cascade gives
\begin{align}
\Delta_2&=bK_p-\tau_mb(K_i+k_aK_p),\\
\Delta_3&=b(K_i+k_aK_p)\Delta_2-bk_aK_i.
\end{align}
For altitude use $a_3=1+\gamma\tau_m$, $a_2=\gamma+b_hK_d$,
$a_1=b_hK_p$ and $a_0=b_hK_i$. The executable analysis reports each
determinant, every pole and the actual step metrics for both gain sets.
These strict tests apply to the local unsaturated models. Ground contact,
mode transitions and a changing camera map are not a single LTI system.

\section{A continuous parameter-interval proof for the conservative candidate}
The grid finds counterexamples. For $\lambda=0.7$ we can additionally prove
local Routh inequalities over a continuous box. Assume
$b\in[0.65b_0,1.35b_0]$, $\tau\in[0.02,0.09]$ s, and for height
$\gamma\in[0.08,0.13]$ s$^{-1}$. These are assumed ranges, not measured
confidence intervals for the real aircraft.

For the angle cascade define $u=K_i+k_aK_p>0$ and
$A(\tau)=K_p-\tau u$. Since $u>0$, $A$ decreases with $\tau$, so
$A\geq A_{\min}=K_p-\tau_{\max}u$. The determinants become
\begin{equation}
\Delta_2=bA,\qquad\Delta_3=b[buA-k_aK_i].
\end{equation}
If $A_{\min}>0$, multiplication by positive $b,u$ preserves inequalities.
It is therefore sufficient to check
\begin{equation}
A_{\min}>0,\qquad b_{\min}uA_{\min}-k_aK_i>0.
\end{equation}
The numerical lower bounds are $0.000257095$ and $0.00776085$ approximately.
All coefficients and Routh first-column entries are consequently positive
throughout that continuous box for this decoupled angle model. Yaw requires
$K_p-\tau_{\max}K_i>0$, with lower bound $0.00462303$.

For height, expand the determinant without skipping terms:
\begin{align}
\Delta_2&=(1+\gamma\tau)(\gamma+bK_d)-\tau bK_p\\
&=\gamma+\gamma^2\tau+bK_d+b\gamma\tau K_d-b\tau K_p.
\end{align}
Since the discarded $\gamma^2\tau$ and $b\gamma\tau K_d$ terms are nonnegative,
and $K_d-\tau_{\max}K_p>0$ for this candidate,
\begin{equation}
\Delta_2\geq\gamma_{\min}+b_{\min}(K_d-\tau_{\max}K_p)
 =D_{2,\rm lb}>0.
\end{equation}
The next determinant obeys
\begin{align}
\Delta_3&=b\left[K_p\Delta_2-K_i(1+\gamma\tau)^2\right]\\
&\geq b\left[K_pD_{2,\rm lb}-K_i(1+\gamma_{\max}\tau_{\max})^2\right].
\end{align}
Here $D_{2,\rm lb}\simeq1.5460874$ and the final bracket is approximately
$0.1354504>0$. This proves strict local height-loop stability over the specified
continuous box. It excludes saturation, nonlinear coupling, estimator dynamics,
contacts, delays and companion guidance. Those exclusions are part of the theorem.

The fast candidate also has an analytic counterexample. For $\lambda=1.3$,
$K_i/K_p=8.6667$ s$^{-1}$. The necessary angle condition $A>0$ implies
$k_a<1/\tau-K_i/K_p$. At $\tau=0.09$ s the right side is
$2.4444$ s$^{-1}$, below the chosen $k_a=3$. Thus this Routh condition fails
regardless of positive plant gain at that lag. Nominal success at 0.04 s
does not establish stability at 0.09 s.

\section{The horizontal position loop is sixth order}
The companion generates $a_c=k_x(x_c-x)-k_v\dot x$ and $\theta_c=a_c/g$
in radians. Use $X/\Theta=g/[s(s+\gamma)]$ and the complete angle cascade
$T_\theta=N_\theta/D_\theta$. Substitution gives
\begin{align}
X&=\frac{T_\theta}{s(s+\gamma)}[k_xX_c-(k_x+k_vs)X],\\
\big[s(s+\gamma)D_\theta+(k_vs+k_x)N_\theta\big]X&=k_xN_\theta X_c,\\
\boxed{\frac{X}{X_c}}&=\boxed{\frac{k_xN_\theta}
 {s(s+\gamma)D_\theta+(k_vs+k_x)N_\theta}}.
\end{align}
This denominator is sixth order. Navigation uses $k_x=1.8$ s$^{-2}$,
$k_v=1.8$ s$^{-1}$; approach/landing uses $k_x=1.2$ and the same $k_v$.
The lateral-roll sign cancels with its guidance-conversion sign, giving the
same local scalar polynomial under the symmetric model. This is upstream
guidance; there are still exactly six flight-controller PID objects.

For $D_x=d_6s^6+d_5s^5+\cdots+d_0$, the specific Hurwitz array is
\begin{equation}
\mathcal H=\begin{bmatrix}
d_5&d_3&d_1&0&0&0\\d_6&d_4&d_2&d_0&0&0\\
0&d_5&d_3&d_1&0&0\\0&d_6&d_4&d_2&d_0&0\\
0&0&d_5&d_3&d_1&0\\0&0&d_6&d_4&d_2&d_0
\end{bmatrix}.
\end{equation}
For $d_6>0$ all six leading principal determinants must be positive. The
executable analysis reports every determinant and pole for each candidate.
Speed limits, changing waypoints and camera-map updates remain nonlinear
effects beyond this local calculation.

\section{Discrete-time stability at the real control interval}
Keep the physical state $x_p=[\phi,p,u_{\rm act}]^{\mathsf T}$ and integral
state $I$. During one held command, obtain the exact linear plant update
$x_{p,k+1}=A_dx_{p,k}+B_du_k$ by solving the linear ODE over 2 ms.
The software uses the current error in its candidate integral:
\begin{align}
e_k&=k_a(\phi_{c,k}-\phi_k)-p_k,\\
I_{k+1}&=I_k+\Delta e_k,\\
u_k&=K_pe_k+K_iI_{k+1}=(K_p+K_i\Delta)e_k+K_iI_k.
\end{align}
Let $C_e=[-k_a,-1,0]$. At zero reference the combined update is therefore
\begin{equation}
\begin{bmatrix}x_{p,k+1}\\I_{k+1}\end{bmatrix}=
\begin{bmatrix}
A_d+B_d(K_p+K_i\Delta)C_e&B_dK_i\\
\Delta C_e&1
\end{bmatrix}
\begin{bmatrix}x_{p,k}\\I_k\end{bmatrix}.
\end{equation}
This explicitly includes the integration ordering used in C. Strict sampled
stability requires every eigenvalue $z_i$ to satisfy $|z_i|<1$. Yaw uses the
same construction without the angle state; altitude uses $C_e=[-1,0,0]$
and velocity feedback $-K_dv$. Estimator dynamics, computation delay and
allocation limits are additional effects not present in this particular test.
The numerical analysis also repeats it with one control sample of input delay
and reports the resulting spectral radius rather than asserting timing safety.

\section{Include the 20 Hz companion with the 500 Hz flight loop}
For one horizontal axis use physical state
$x_p=[\theta_d,p_d,u_{\rm act},x,v]^{\mathsf T}$. Its hover plant is
\begin{equation}
\dot x_p=\begin{bmatrix}
0&1&0&0&0\\0&0&b_x&0&0\\0&0&-1/\tau_m&0&0\\
0&0&0&0&1\\g\pi/180&0&0&0&-\gamma
\end{bmatrix}x_p+
\begin{bmatrix}0\\0\\1/\tau_m\\0\\0\end{bmatrix}u.
\end{equation}
The $g\pi/180$ entry explicitly converts the degree-valued angle into the
radian-valued small-angle force law. Obtain $A_d,B_d$ for one held 2 ms
command, append the PI integral, and apply the same current-error integral
ordering as above. With $k=K_p+K_i\Delta$, $C_e=[-k_a,-1,0,0,0]$,
\begin{equation}
F=\begin{bmatrix}A_d+B_dkC_e&B_dK_i\\\Delta C_e&1\end{bmatrix},\qquad
G=\begin{bmatrix}B_dkk_a\\\Delta k_a\end{bmatrix}.
\end{equation}
Thus $x_{k+1}=Fx_k+G\theta_c$ for the six-state augmented system.
The companion holds $\theta_c=C_gx_k$ for $N=25$ flight updates, where
\begin{equation}
C_g=-\frac{180}{\pi g}[0,0,0,k_x,k_v,0].
\end{equation}
Substitute the recurrence a second time:
$x_{k+2}=F^2x_k+(F+I)G\theta_c$. A third substitution gives
$F^3x_k+(F^2+F+I)G\theta_c$. Continuing the same physical recurrence for
25 held steps gives
\begin{equation}
x_{k+N}=F^Nx_k+\left(\sum_{j=0}^{N-1}F^jG\right)C_gx_k=Mx_k.
\end{equation}
The numerical analysis checks every eigenvalue of $M$. With one complete
guidance-period delay, append the old held command and check
\begin{equation}
M_{\rm delayed}=\begin{bmatrix}
F^N&\sum_{j=0}^{N-1}F^jG\\C_g&0
\end{bmatrix}.
\end{equation}
This is an explicit 500 Hz/20 Hz sampled-system calculation. It still assumes
exact state feedback, linear hover dynamics and scheduled updates. It is not
a proof for the Kalman estimator, delayed perception, saturation, or Windows
meeting every real-time deadline.

\section{Frequency-domain margins and uncertainty}
For the rate loop, open-loop $L_r=C G=b(K_ps+K_i)/[s^2(\tau_ms+1)]$.
At the gain crossover $|L(j\omega_c)|=1$, phase margin is
$180^\circ+\arg L(j\omega_c)$. At a $-180^\circ$ phase crossover,
gain margin is $1/|L(j\omega_\pi)|$. These are frequency-domain local robustness
measures, not obstacle-clearance margins. A pure delay adds phase
$-\omega T_d$; the rough additional-delay allowance from a positive phase
margin is $T_d\simeq\mathrm{PM}_{\rm rad}/\omega_c$ for the crossing being
considered. Multiple crossovers and unstable open-loop cases need care.

Plant uncertainty enters directly: $b_x\propto aT_{\max}/J_x$ and
$b_h\propto T_{\max}/m$. The executable analysis varies these gains and
$\tau_m$ over a declared finite grid and counts stable/unstable cases for the
complete characteristic polynomials. A grid with no failures is evidence
only at those points; it is not a continuous robust-stability proof. Battery
sag, thrust asymmetry, flexible arms and estimator bias remain unmodelled
physical risks until measured. Time-domain runs then test saturation, takeoff,
landing and perception-driven commands beyond the linear analysis.

\chapter{Companion missions, boost and actuator authority}
\section{Keep one flight-control mode; move mission sequencing upstream}
The flight core still receives roll/pitch angles, yaw rate and centered
vertical-speed throttle. The companion selects mission phases: ground,
spool-up in simulation, takeoff, hover, navigate, approach/hold, descent and
landed. A mission phase does not introduce a new flight PID. The companion
converts a requested height into bounded vertical speed; the existing altitude
PID integrates that speed into its internal target. Start and finish points
are explicit local coordinates. The start must agree with the measured launch
position; selecting it must never teleport the aircraft state.

Ground contact supplies an upward normal reaction that balances weight before
lift-off. At $T_z>mg$ that unilateral reaction vanishes and upward acceleration
begins. Landing is the reverse transition with low descent speed and contact
confirmation; simply forcing a trajectory through $h=0$ is not a contact model.
The 3D simulator handles contact; real landed detection remains a separate
sensor/state-machine problem. A simulation-only motor gate cannot be presented
as an implemented physical arming system.

\section{What arming and RT boost mean}
Arming is a deliberate permission to produce physical motor outputs. It is not
the same as compiling a control function or connecting a battery. Physical
arming stays disabled as requested. Takeoff/landing button gestures are mission
requests; a future authorized safety supervisor must accept them before motors
can be enabled. The simulated actuator gate is explicitly isolated from hardware.

RT is already present in the radio frame as a normalized trigger. Its proposed
boost increases the maximum commanded vertical speed from 2 to 3 m/s and the
nominal collective ceiling from 0.65 to 0.80 as RT moves from zero to one.
The internal boost allowance rises over at least 0.25 s and decays over at least
0.10 s for a full-scale change. RT alone does not command a climb at centered
vertical stick. The companion simulation inhibits boost during approach/landing.
Roll/pitch leveling and yaw-rate control remain active. It is not acrobatic
rate mode and not an unconditional full-throttle override. The boost ceiling
has less correction reserve than normal mode; maximum climb requires measured
thrust/energy reserve and later hardware tuning.

\section{Prove the motor-headroom allocation interval}
Let desired motor values before clipping be $u_i=c+d_i$, where $d_i$ are the
four differential roll/pitch/yaw corrections. For each motor to satisfy
$0\leq u_i\leq1$, rearrange to $-d_i\leq c\leq1-d_i$.
Intersect all four intervals:
\begin{equation}
-\min_i d_i\leq c\leq1-\max_i d_i.
\end{equation}
This interval exists exactly when $\max d_i-\min d_i\leq1$. If the requested
differential span exceeds one, uniformly scale all $d_i$ by
$\alpha=1/(\max d_i-\min d_i)$; otherwise $\alpha=1$. Clamp collective to
the resulting feasible interval and calculate $u_i=c+\alpha d_i$.
This preserves differential-command ratios, unlike clipping each motor
independently. It cannot create missing thrust or guarantee full simultaneous
torque authority. A future allocator should feed allocation residuals back into
anti-windup; the current PID anti-windup covers each loop's own output limit,
not the complete coupled allocation constraint.
